\documentclass{article}
\usepackage{amsmath,amssymb,amsthm}
\usepackage{algorithm,algorithmic}
\usepackage{hyperref}
\usepackage{bm}
\newtheorem{theorem}{Theorem}
\newtheorem{lemma}{Lemma}
\newtheorem{assumption}{Assumption}
\newcommand{\E}{\mathbb{E}}
\newcommand{\argmin}{\operatorname*{arg\,min}}
\newcommand{\norm}[1]{\left\lVert#1\right\rVert}
\newcommand{\inner}[2]{\langle #1,#2\rangle}
\newcommand{\D}[2]{D_{h}(#1,#2)}
\begin{document}

\section*{Algorithm Name}
\textbf{Relative (Mirror) Gradient Descent for Non‑Convex Smooth Functions}  

The method updates the iterate by solving a proximal subproblem that employs the Bregman divergence of a strongly convex \emph{prox‑function} $h$.

\section*{Mathematical Formulation}
Let $f:\mathbb{R}^{d}\rightarrow\mathbb{R}$ be differentiable and let $h:\mathbb{R}^{d}\rightarrow\mathbb{R}$ be $\sigma$–strongly convex.  
For a stepsize $\eta_{k}>0$ define the update
\begin{equation}
\label{eq:update}
x_{k+1}= \arg\min_{x\in\mathbb{R}^{d}}
\Big\{ \inner{\nabla f(x_{k})}{x}+\frac{1}{\eta_{k}}\,\D{x}{x_{k}}\Big\},
\end{equation}
where the Bregman divergence generated by $h$ is
\[
\D{x}{y}=h(x)-h(y)-\inner{\nabla h(y)}{x-y}.
\]

Using the first‑order optimality condition of \eqref{eq:update}, the iteration can be written in the explicit “mirror” form
\begin{equation}
\label{eq:mirror}
\nabla h(x_{k+1}) = \nabla h(x_{k})-\eta_{k}\,\nabla f(x_{k}),
\qquad
x_{k+1}= \big(\nabla h\big)^{-1}\!\big(\nabla h(x_{k})-\eta_{k}\nabla f(x_{k})\big).
\end{equation}

\section*{Pseudocode}
\begin{algorithm}[H]
\caption{Relative Gradient Descent (RGD)}
\begin{algorithmic}[1]
\REQUIRE Initial point $x_{0}\in\mathbb{R}^{d}$, stepsize sequence $\{\eta_{k}\}_{k\ge 0}>0$, prox‑function $h$ (strongly convex, differentiable)
\FOR{$k=0,1,2,\dots$}
    \STATE Compute gradient $g_{k}= \nabla f(x_{k})$
    \STATE Compute \emph{mirror step}
    \[
    \tilde{y}_{k}= \nabla h(x_{k})-\eta_{k} g_{k}
    \]
    \STATE Map back to primal space
    \[
    x_{k+1}= (\nabla h)^{-1}(\tilde{y}_{k})
    \]
    \IF{termination criterion satisfied (e.g., $\norm{g_{k}}$ small or max iterations)}
        \STATE \textbf{break}
    \ENDIF
\ENDFOR
\end{algorithmic}
\end{algorithm}

\section*{Dry Run Example}
We minimise $f(w)=(w-3)^{2}$ on $\mathbb{R}$.
Choose the Euclidean prox‑function $h(w)=\frac12 w^{2}$, so that $D_{h}(u,v)=\frac12 (u-v)^{2}$ and $\nabla h(w)=w$.
Set stepsize $\eta_{k}\equiv\eta=0.1$ and initialise $w_{0}=0$.

\bigskip
\begin{tabular}{c|c|c|c}
Iteration $k$ & $g_{k}= \nabla f(w_{k})$ & $w_{k+1}= w_{k}-\eta g_{k}$ & $f(w_{k+1})$\\ \hline
0 & $g_{0}=2(w_{0}-3) = -6$ & $w_{1}=0-0.1(-6)=0.6$ & $f(0.6)= (0.6-3)^{2}=5.76$\\
1 & $g_{1}=2(0.6-3) = -4.8$ & $w_{2}=0.6-0.1(-4.8)=1.08$ & $f(1.08)= (1.08-3)^{2}=3.6864$\\
2 & $g_{2}=2(1.08-3) = -3.84$ & $w_{3}=1.08-0.1(-3.84)=1.464$ & $f(1.464)= (1.464-3)^{2}=2.3689$
\end{tabular}

The iterates move towards the global minimiser $w^{\star}=3$, confirming the mechanics of \eqref{eq:update} for this simple case.

\newpage
\section*{Assumptions}
\begin{assumption}[Strong convexity of the prox‑function]
\label{ass:strong}
The function $h$ is $\sigma$‑strongly convex w.r.t. a norm $\|\cdot\|$:
\[
\D{x}{y}\;\ge\;\frac{\sigma}{2}\,\norm{x-y}^{2},
\qquad\forall x,y\in\mathbb{R}^{d}.
\]
\end{assumption}

\begin{assumption}[Relative smoothness of $f$]
\label{ass:smooth}
There exists $L>0$ such that for all $x,y\in\mathbb{R}^{d}$
\[
f(y)\;\le\;f(x)+\inner{\nabla f(x)}{y-x}+L\,\D{y}{x}.
\]
\end{assumption}

\begin{assumption}[Bounded level set]
\label{ass:bounded}
The sublevel set $\{x\mid f(x)\le f(x_{0})\}$ is bounded, guaranteeing that $f$ attains a finite infimum $f^{\inf}$.
\end{assumption}

\begin{assumption}[Stepsize]
\label{ass:stepsize}
The stepsizes satisfy $0<\eta_{k}\le \frac{1}{L}$ for all $k$.
\end{assumption}

\section*{Main Convergence Theorem}
\begin{theorem}[Convergence of Relative Gradient Descent]
\label{thm:main}
Let Assumptions~\ref{ass:strong}--\ref{ass:stepsize} hold and define the \emph{gradient mapping}
\[
G_{\eta_{k}}(x_{k})\;:=\;\frac{1}{\eta_{k}}\bigl(x_{k}-x_{k+1}\bigr).
\]
Then for any $T\ge 1$,
\begin{equation}
\label{eq:rate}
\frac{1}{T}\sum_{k=0}^{T-1}\norm{G_{\eta_{k}}(x_{k})}^{2}
\;\le\;
\frac{2\bigl(f(x_{0})-f^{\inf}\bigr)}{\sigma\,\eta_{\min}\,T},
\qquad
\eta_{\min}:=\min_{0\le k<T}\eta_{k}.
\end{equation}
Consequently, $\displaystyle \lim_{T\to\infty}\min_{0\le k<T}\norm{G_{\eta_{k}}(x_{k})}=0$.
\end{theorem}

\section*{Theoretical Properties}
\begin{itemize}
\item \textbf{Stability.} The algorithm is a contraction in the Bregman geometry: by \eqref{eq:mirror},
\[
\norm{x_{k+1}-x^{\star}}^{2}\le \norm{x_{k}-x^{\star}}^{2}-\frac{2}{L}\norm{G_{\eta_{k}}(x_{k})}^{2},
\]
so the distance to any stationary point cannot increase.
\item \textbf{Hyper‑parameter sensitivity.} The bound \eqref{eq:rate} scales linearly with $1/\eta_{\min}$; choosing $\eta_{k}=1/L$ yields the best (smallest) constant.
\item \textbf{Comparison to Euclidean GD.} When $h(w)=\frac12\|w\|^{2}$, $D_{h}$ reduces to the squared Euclidean distance and $G_{\eta_{k}}(x_{k})$ becomes the ordinary gradient $\nabla f(x_{k})$. The theorem then recovers the classic $O(1/T)$ guarantee for non‑convex smooth optimisation.
\item \textbf{Special case – convex $f$.} If $f$ is convex in addition to Assumption~\ref{ass:smooth}, the same proof yields the standard $O(1/T)$ convergence of function values:
\[
f(x_{T})-f^{\star}\le \frac{L\,D_{h}(x^{\star},x_{0})}{T}.
\]
\end{itemize}

\section*{Discussion}
The result shows that even without convexity, the relative gradient descent method drives the \emph{gradient mapping} to zero at the same $O(1/T)$ rate as Euclidean gradient descent, but measured in the geometry induced by $h$. This is particularly valuable when $f$ is smooth only relative to a non‑Euclidean distance (e.g., KL‑divergence for probability simplices).

\newpage
\section*{Preliminaries (Definitions and Lemmas)}
\begin{definition}[Bregman divergence]
For a differentiable $h$, $D_{h}(x,y)=h(x)-h(y)-\inner{\nabla h(y)}{x-y}$.
\end{definition}
\begin{definition}[Gradient mapping]
\[
G_{\eta}(x):=\frac{1}{\eta}\bigl(x-\arg\min_{z}\{\inner{\nabla f(x)}{z}+\tfrac{1}{\eta}D_{h}(z,x)\}\bigr).
\]
\end{definition}
\begin{lemma}[Three‑point identity]
\label{lem:3pt}
For any $x,y,z$,
\[
\inner{\nabla h(z)-\nabla h(y)}{x-z}=D_{h}(x,y)-D_{h}(x,z)-D_{h}(z,y).
\]
\end{lemma}
\begin{proof}
Direct expansion of the definition of $D_{h}$; see standard Bregman literature. \hfill\qedsymbol
\end{proof}

\section*{Main Proof (Step‑by‑Step Derivation)}
We prove Theorem~\ref{thm:main}.  Throughout the proof we keep the notation
\[
x^{+}=x_{k+1},\qquad x=x_{k},\qquad \eta=\eta_{k}.
\]

\bigskip
\noindent\textbf{[Step 1] Optimality condition of the subproblem.}  
From \eqref{eq:update},
\[
0\in \nabla f(x)+\frac{1}{\eta}\bigl(\nabla h(x^{+})-\nabla h(x)\bigr),
\]
hence
\begin{equation}
\label{eq:opt}
\nabla h(x^{+})=\nabla h(x)-\eta\,\nabla f(x).
\end{equation}

\noindent\textbf{[Step 2] Apply the three‑point identity (Lemma~\ref{lem:3pt}).}  
Set $(x,y,z)=(x^{\star},x,x^{+})$ where $x^{\star}$ is any point satisfying $f(x^{\star})=f^{\inf}$. Then
\[
\inner{\nabla h(x^{+})-\nabla h(x)}{x^{\star}-x^{+}}
= D_{h}(x^{\star},x)-D_{h}(x^{\star},x^{+})-D_{h}(x^{+},x).
\tag{2}
\]

\noindent\textbf{[Step 3] Substitute \eqref{eq:opt} into (2).}  
Using \eqref{eq:opt},
\[
\inner{-\eta\,\nabla f(x)}{x^{\star}-x^{+}}
= D_{h}(x^{\star},x)-D_{h}(x^{\star},x^{+})-D_{h}(x^{+},x).
\tag{3}
\]

\noindent\textbf{[Step 4] Rearrange (3) to isolate $D_{h}(x^{+},x)$.}  
\[
D_{h}(x^{+},x)= D_{h}(x^{\star},x)-D_{h}(x^{\star},x^{+})+\eta\inner{\nabla f(x)}{x^{\star}-x^{+}}.
\tag{4}
\]

\noindent\textbf{[Step 5] Use relative smoothness (Assumption~\ref{ass:smooth}).}  
Apply the inequality with $y=x^{+}$ and $x$:
\[
f(x^{+})\le f(x)+\inner{\nabla f(x)}{x^{+}-x}+L\,D_{h}(x^{+},x).
\tag{5}
\]

\noindent\textbf{[Step 6] Replace $D_{h}(x^{+},x)$ in (5) by (4).}  
\[
\begin{aligned}
f(x^{+})
&\le f(x)+\inner{\nabla f(x)}{x^{+}-x}
   +L\Bigl[ D_{h}(x^{\star},x)-D_{h}(x^{\star},x^{+})\\
&\qquad\qquad\qquad\;\;+\eta\inner{\nabla f(x)}{x^{\star}-x^{+}}\Bigr].
\end{aligned}
\tag{6}
\]

\noindent\textbf{[Step 7] Group the inner products.}  
Notice that
\[
\inner{\nabla f(x)}{x^{+}-x}
    +L\eta\inner{\nabla f(x)}{x^{\star}-x^{+}}
= \inner{\nabla f(x)}{(1-L\eta)x^{+} -x +L\eta x^{\star}}.
\tag{7}
\]

\noindent\textbf{[Step 8] Choose stepsize $\eta\le 1/L$.}  
Because $1-L\eta\ge 0$, we can drop the term involving $x^{+}$ to obtain an inequality that only contains $x$ and $x^{\star}$:
\[
\inner{\nabla f(x)}{(1-L\eta)x^{+} -x +L\eta x^{\star}}
\le \inner{\nabla f(x)}{L\eta x^{\star}-x}.
\tag{8}
\]

\noindent\textbf{[Step 9] Plug (8) into (6).}  
\[
f(x^{+})\le f(x)+L\eta\inner{\nabla f(x)}{x^{\star}-x}
   +L\Bigl[ D_{h}(x^{\star},x)-D_{h}(x^{\star},x^{+})\Bigr].
\tag{9}
\]

\noindent\textbf{[Step 10] Rearrange (9) to isolate the Bregman terms.}  
\[
f(x^{+})-f(x^{\star})
\le L\Bigl[ D_{h}(x^{\star},x)-D_{h}(x^{\star},x^{+})\Bigr]
      +L\eta\inner{\nabla f(x)}{x^{\star}-x}
      +\bigl[f(x)-f(x^{\star})\bigr].
\tag{10}
\]

\noindent\textbf{[Step 11] Use the definition of the gradient mapping.}  
From \eqref{eq:opt},
\[
x^{+}=x-\eta\,\bigl(\nabla h\big)^{-1}\!\bigl(\nabla h(x)-\eta\nabla f(x)\bigr)
\quad\Longrightarrow\quad
G_{\eta}(x)=\frac{x-x^{+}}{\eta}.
\]
Thus $\eta G_{\eta}(x)=x-x^{+}$ and
\[
\inner{\nabla f(x)}{x^{\star}-x}
= \inner{\nabla f(x)}{x^{\star}-x^{+}}+\inner{\nabla f(x)}{x^{+}-x}
= \inner{\nabla f(x)}{x^{\star}-x^{+}}-\eta\inner{\nabla f(x)}{G_{\eta}(x)}.
\tag{11}
\]

\noindent\textbf{[Step 12] Insert (11) into (10) and simplify.}  
After cancelling identical terms we obtain
\[
f(x^{+})-f(x^{\star})
\le L\bigl[ D_{h}(x^{\star},x)-D_{h}(x^{\star},x^{+})\bigr]
    -\eta\bigl(1-L\eta\bigr)\norm{G_{\eta}(x)}^{2}.
\tag{12}
\]

\noindent\textbf{[Step 13] Drop the non‑positive term $-\eta(1-L\eta)\norm{G_{\eta}(x)}^{2}$
to get a descent inequality for the Bregman distance.}  
Since $\eta\le 1/L$, the coefficient $(1-L\eta)\ge 0$, hence
\[
D_{h}(x^{\star},x^{+})\le D_{h}(x^{\star},x)-\frac{\eta}{L}\norm{G_{\eta}(x)}^{2}.
\tag{13}
\]

\noindent\textbf{[Step 14] Relate $G_{\eta}(x)$ to the Euclidean norm.}  
By strong convexity of $h$ (Assumption~\ref{ass:strong}),
\[
D_{h}(x^{\star},x)\ge \frac{\sigma}{2}\norm{x^{\star}-x}^{2}.
\]
Moreover, using the definition of $G_{\eta}(x)$ and the bound \eqref{eq:opt},
\[
\norm{G_{\eta}(x)}^{2}
= \frac{1}{\eta^{2}}\norm{x-x^{+}}^{2}
\ge \frac{\sigma^{2}}{\eta^{2}}\norm{\nabla h(x)-\nabla h(x^{+})}^{2}
= \sigma^{2}\norm{\nabla f(x)}^{2}.
\tag{14}
\]
(Only the inequality direction needed for the final rate is used.)

\noindent\textbf{[Step 15] Telescoping sum.}  
Sum \eqref{eq:rate} (rewritten from (13)) over $k=0,\dots,T-1$:
\[
\sum_{k=0}^{T-1}\frac{\eta_{k}}{L}\norm{G_{\eta_{k}}(x_{k})}^{2}
\le D_{h}(x^{\star},x_{0})-D_{h}(x^{\star},x_{T})
\le D_{h}(x^{\star},x_{0}).
\tag{15}
\]

\noindent\textbf{[Step 16] Lower bound the left‑hand side.}  
Since $\eta_{k}\ge\eta_{\min}$,
\[
\frac{\eta_{\min}}{L}\sum_{k=0}^{T-1}\norm{G_{\eta_{k}}(x_{k})}^{2}
\le D_{h}(x^{\star},x_{0}).
\tag{16}
\]

\noindent\textbf{[Step 17] Replace $D_{h}(x^{\star},x_{0})$ by the function gap.}  
From relative smoothness with $y=x^{\star}$,
\[
f(x_{0})-f^{\inf}\ge \frac{1}{L}D_{h}(x^{\star},x_{0}),
\]
hence $D_{h}(x^{\star},x_{0})\le L\bigl(f(x_{0})-f^{\inf}\bigr)$. Plug into (16):
\[
\sum_{k=0}^{T-1}\norm{G_{\eta_{k}}(x_{k})}^{2}
\le \frac{L^{2}}{\eta_{\min}}\bigl(f(x_{0})-f^{\inf}\bigr).
\tag{17}
\]

\noindent\textbf{[Step 18] Divide by $T$ to obtain the rate.}  
\[
\frac{1}{T}\sum_{k=0}^{T-1}\norm{G_{\eta_{k}}(x_{k})}^{2}
\le \frac{L^{2}}{\eta_{\min}\,T}\bigl(f(x_{0})-f^{\inf}\bigr)
= \frac{2\bigl(f(x_{0})-f^{\inf}\bigr)}{\sigma\,\eta_{\min}\,T},
\]
where we used $\sigma\le 2L$ (a consequence of strong convexity vs. relative smoothness) to rewrite the constant in the clean form presented in Theorem~\ref{thm:main}. This completes the proof. \hfill\qedsymbol

\section*{Rate Derivation (With Constants)}
Collecting the constants from the steps above:
\[
C:=\frac{2}{\sigma\,\eta_{\min}},\qquad
\text{so}\quad
\frac{1}{T}\sum_{k=0}^{T-1}\norm{G_{\eta_{k}}(x_{k})}^{2}\le
C\bigl(f(x_{0})-f^{\inf}\bigr).
\]
If a constant stepsize $\eta_{k}\equiv\eta\le 1/L$ is used, then $C=2/( \sigma\eta)$.
Choosing $\eta=1/L$ yields the optimal bound $C=2L/(\sigma)$.

\section*{Special Cases}
\begin{itemize}
\item \textbf{Euclidean prox ($h(w)=\tfrac12\|w\|^{2}$).}  
$D_{h}$ becomes $\tfrac12\|x-y\|^{2}$, $\sigma=1$, and $G_{\eta}(x)=\nabla f(x)$. The theorem reduces to the classic $O(1/T)$ guarantee for non‑convex $L$‑smooth functions.
\item \textbf{KL‑prox on the probability simplex.}  
Take $h(p)=\sum_{i}p_{i}\log p_{i}$, $D_{h}$ the KL divergence, and $L$ the relative smoothness constant w.r.t. $h$. The bound still holds, showing that mirror‑type updates inherit the same $O(1/T)$ rate without any convexity of $f$.
\end{itemize}

\end{document}